\documentclass[10pt,a4paper]{amsart}
\usepackage{amsmath,amssymb,amsthm,graphicx}
% Retained minimal preamble for the Stacks Project chapter sites-cohomology.tex (Cohomology on Sites).
% Verbatim excerpt of the chapter preamble: Xy-pic setup (native preamble lines 14-18) and the native
% theorem declarations with the mathematical macros (native preamble lines 154-217). Copied mechanically;
% nothing added, removed or redefined. The wrapper class and the amsmath/amssymb/amsthm/graphicx base
% packages are supplied by the assembler, so the class/externaldocument part of the native preamble is omitted.
% For commutative diagrams we use Xy-pic
\usepackage[all]{xy}

% We use 2cell for 2-commutative diagrams.
\xyoption{2cell}
% Theorem environments.
%
\theoremstyle{plain}
\newtheorem{theorem}[subsection]{Theorem}
\newtheorem{proposition}[subsection]{Proposition}
\newtheorem{lemma}[subsection]{Lemma}

\theoremstyle{definition}
\newtheorem{definition}[subsection]{Definition}
\newtheorem{example}[subsection]{Example}
\newtheorem{exercise}[subsection]{Exercise}
\newtheorem{situation}[subsection]{Situation}

\theoremstyle{remark}
\newtheorem{remark}[subsection]{Remark}
\newtheorem{remarks}[subsection]{Remarks}

\numberwithin{equation}{subsection}

% Macros
%
\def\lim{\mathop{\mathrm{lim}}\nolimits}
\def\colim{\mathop{\mathrm{colim}}\nolimits}
\def\Spec{\mathop{\mathrm{Spec}}}
\def\Hom{\mathop{\mathrm{Hom}}\nolimits}
\def\Ext{\mathop{\mathrm{Ext}}\nolimits}
\def\SheafHom{\mathop{\mathcal{H}\!\mathit{om}}\nolimits}
\def\SheafExt{\mathop{\mathcal{E}\!\mathit{xt}}\nolimits}
\def\Sch{\mathit{Sch}}
\def\Mor{\mathop{\mathrm{Mor}}\nolimits}
\def\Ob{\mathop{\mathrm{Ob}}\nolimits}
\def\Sh{\mathop{\mathit{Sh}}\nolimits}
\def\NL{\mathop{N\!L}\nolimits}
\def\CH{\mathop{\mathrm{CH}}\nolimits}
\def\proetale{{pro\text{-}\acute{e}tale}}
\def\etale{{\acute{e}tale}}
\def\QCoh{\mathit{QCoh}}
\def\Ker{\mathop{\mathrm{Ker}}}
\def\Im{\mathop{\mathrm{Im}}}
\def\Coker{\mathop{\mathrm{Coker}}}
\def\Coim{\mathop{\mathrm{Coim}}}

% Boxtimes
%
\DeclareMathSymbol{\boxtimes}{\mathbin}{AMSa}{"02}

%
% Macros for moduli stacks/spaces
%
\def\QCohstack{\mathcal{QC}\!\mathit{oh}}
\def\Cohstack{\mathcal{C}\!\mathit{oh}}
\def\Spacesstack{\mathcal{S}\!\mathit{paces}}
\def\Quotfunctor{\mathrm{Quot}}
\def\Hilbfunctor{\mathrm{Hilb}}
\def\Curvesstack{\mathcal{C}\!\mathit{urves}}
\def\Polarizedstack{\mathcal{P}\!\mathit{olarized}}
\def\Complexesstack{\mathcal{C}\!\mathit{omplexes}}
% \Pic is the operator that assigns to X its picard group, usage \Pic(X)
% \Picardstack_{X/B} denotes the Picard stack of X over B
% \Picardfunctor_{X/B} denotes the Picard functor of X over B
\def\Pic{\mathop{\mathrm{Pic}}\nolimits}
\def\Picardstack{\mathcal{P}\!\mathit{ic}}
\def\Picardfunctor{\mathrm{Pic}}
\def\Deformationcategory{\mathcal{D}\!\mathit{ef}}
\usepackage{cleveref}
\crefname{originalequation}{original equation}{original equations}
\Crefname{originalequation}{Original equation}{Original equations}
\makeatletter
\newlabel{equation-projection-formula-map}{{0B56}{source}{}{sourcepointer-1}{}}
\newlabel{equation-projection-formula-map@cref}{{[originalequation][0][]0B56}{[1][0][]source}{}{}{}}
\makeatother
\title{Projection formula for a perfect coefficient object}
\author{The Stacks Project Authors}
\date{stacks-project:sites-cohomology.tex:local-snapshot-20261011; GNU FDL 1.2 or later}
\begin{document}
\maketitle
\noindent\textit{Editorial note.} All delivered mathematical text is the author's own, copied from the pinned Stacks Project chapter file sites-cohomology.tex; no statement or proof step was rewritten, reordered or removed.\par\medskip
\noindent\textit{Editorial note.} The proof is the author's complete reduction: the four steps and the trivial case are all stated in the counted span, including the additivity step and the truncation step.\par\medskip
\noindent\textit{Editorial note.} The single labelled reference occurring inside the counted spans is carried in assembly.external\_references as a non-load-bearing navigation pointer to the numbered equation of the same chapter that defines the map being checked.\par\medskip
\noindent\textit{Editorial note.} Cross-chapter restatement: the same label occurs in chapter cohomology for a ringed space and a finite locally free coefficient, which is the special case of this perfect-coefficient statement, so per the lane-s4 rule only this more general occurrence is counted here. The same label also occurs in chapter chow for a different proposition (the Gysin compatibility of proper smooth morphisms), which is a name collision only, not a restatement.\par\medskip
\noindent\textit{Editorial note.} The preamble is the chapter's own retained minimal preamble (sources/stacks-sites-cohomology/retained-preamble.tex), a verbatim excerpt of the native preamble: the Xy-pic setup and the native theorem declarations with the mathematical macros.\par\medskip
\noindent\textit{Allowed standard prerequisites (editorial).}\begin{itemize}
\item Perfect object of D(O) on a ringed site. Same chapter, the definition of perfect objects: locally on the site represented by a strictly perfect complex. Supplies the hypothesis on K and the reduction to a strictly perfect representative used in the proof.
\item Strictly perfect complex and finite free modules. Same chapter, the description of a strictly perfect complex as a bounded complex of direct summands of finite free modules. Supplies the second reduction, from a strictly perfect complex to a finite complex of finite free modules.
\item Derived pullback, derived pushforward and the derived tensor product of modules on a ringed topos. Same chapter, the functors Lf\textasciicircum{}*, Rf\_* and tensor\textasciicircum{}L appearing on both sides of the asserted identity. Supplies the two sides whose equality is asserted and the trivial case K = O\_D[n].
\end{itemize}\medskip
\ifdefined\hypertarget\hypertarget{sourcepointer-1}{}\fi
\noindent\textit{Source reference.} Original equation 0B56: sources/stacks-sites-cohomology/pinned-original-source.tex, source lines 13666--13666. Same-chapter navigation pointer invoked at proof line 13696, where the author checks that the map displayed in that numbered equation is an isomorphism. The pointer only resolves the reference for the reader; the equation itself is not reproduced inside the counted span, and no bibliography citation is involved.\par\medskip
\begin{lemma}
\label{lemma-projection-formula}
Let $f : (\Sh(\mathcal{C}), \mathcal{O}_\mathcal{C}) \to
(\Sh(\mathcal{D}), \mathcal{O}_\mathcal{D})$ be a morphism of ringed topoi.
Let $E \in D(\mathcal{O}_\mathcal{C})$ and $K \in D(\mathcal{O}_\mathcal{D})$.
If $K$ is perfect, then
$$
Rf_*E \otimes^\mathbf{L}_{\mathcal{O}_\mathcal{D}} K =
Rf_*(E \otimes^\mathbf{L}_{\mathcal{O}_\mathcal{C}} Lf^*K)
$$
in $D(\mathcal{O}_\mathcal{D})$.
\end{lemma}

\begin{proof}
To check (\ref{equation-projection-formula-map}) is an isomorphism
we may work locally on $\mathcal{D}$, i.e.,
for any object $V$ of $\mathcal{D}$ we have to find a covering $\{V_j \to V\}$
such that the map restricts to an isomorphism on $V_j$. By definition
of perfect objects, this means we may assume $K$ is represented by
a strictly perfect complex of $\mathcal{O}_\mathcal{D}$-modules.
Note that, completely generally, the statement is true for
$K = K_1 \oplus K_2$, if and only if the statement is true for
$K_1$ and $K_2$. Hence we may assume $K$ is a finite
complex of finite free $\mathcal{O}_\mathcal{D}$-modules.
In this case a simple argument involving stupid truncations reduces
the statement to the case where $K$ is represented by a finite
free $\mathcal{O}_\mathcal{D}$-module. Since the statement is invariant
under finite direct summands in the $K$ variable, we conclude
it suffices to prove it for $K = \mathcal{O}_\mathcal{D}[n]$
in which case it is trivial.
\end{proof}

\end{document}
